\title{FlashAttention-3: Fast and Accurate Attention with Asynchrony and Low-precision}

\begin{abstract}
Attention mechanisms serve as a crucial component within the Transformer framework, yet they represent a computational bottleneck for both large language models and scenarios requiring extended contexts. The \textsc{FlashAttention} algorithm previously introduced optimizations for GPU attention by significantly reducing memory traffic. Nonetheless, previous iterations have not fully utilized the capabilities of the latest hardware generations; for example, \textsc{FlashAttention-2} is able to attain merely 35\% hardware utilization on the H100 GPU. In our work, we introduce three core strategies aimed at accelerating attention on Hopper GPUs: harnessing the asynchronous features of Tensor Cores and TMA to (1) orchestrate a simultaneous overlap of data transfer and computation through warp specialization, (2) interleave matrix multiplication with softmax operations at the block level, and (3) employ block-level quantization along with non-coherent processing, benefiting from native FP8 precision support in hardware. Our proposed approach, \textsc{FlashAttention-3}, achieves between 1.5$\times$ and 2.0$\times$ speedup on H100 GPUs, attaining up to 740 TFLOPs/s with FP16 (corresponding to 75\% utilization), and nearly 1.2 PFLOPs/s using FP8. We further demonstrate that FP8 \textsc{FlashAttention-3} produces 2.6$\times$ smaller numerical error compared with a standard FP8 attention method.
\end{abstract}

\section{Introduction}
\label{sec:intro}

Within the Transformer architecture~\citep{vaswani2017attention}, the primary computational constraint arises from the attention mechanism. Specifically, self-attention requires computation of attention scores between queries and keys, resulting in quadratic complexity with respect to sequence length. Extending attention to accommodate longer sequences presents opportunities for enhanced capabilities—such as reasoning over multiple extensive documents~\citep{guo2021longt5,shaham2022scrolls,peng2023yarn} or large codebases~\citep{roziere2023code, li2023starcoder}, as well as supporting additional modalities like high-resolution images~\citep{chen2022scaling}, audio~\citep{gulati2020conformer}, and video~\citep{ho2022video}. Further, it can enable new types of applications including user interaction involving extended histories~\citep{sun2019bert4rec} and agent-based workflows with extended horizons~\citep{yao2022react}. As a result, there has been considerable interest in improving the efficiency of attention for long-context scenarios, through various approaches such as approximating attention~\citep{katharopoulos2020transformers,choromanski2020rethinking, tay2020efficient}, optimizing implementation in software (\citep{rabe2021self,dao2022flashattention,kwon2023efficient}), and even developing alternative architectures~\citep{peng2023rwkv,sun2023retentive,gu2023mamba}.

This study extends the research by \citet{dao2022flashattention}, who devised exact-attention algorithms that leverage specific GPU execution paradigms and architectural features during algorithm design. In \citep{dao2022flashattention}, Dao and collaborators presented \textsc{FlashAttention}, introducing an innovative tiling approach to attention computation, which consolidates all attention operations into a single GPU kernel. This strategy eradicates the necessity for intermediate data transfers to and from slower global memory, leading to more efficient parallelism.

Subsequently, \citet{dao2023flashattention2} reformulated the approach as \textsc{FlashAttention-2}, allowing for additional parallelism across the sequence length dimension by executing the forward pass's innermost loop over sub-blocks of the key and value matrices. This redesign yielded improved workload distribution and GPU occupancy. Despite these enhancements, we find that \textsc{FlashAttention-2} still demonstrates suboptimal hardware utilization on the latest GPUs compared to state-of-the-art matrix-multiplication (GEMM) kernels—achieving only 35\% efficiency versus the 80-90\% range reported on the Hopper H100 GPU. One contributing factor may be the absence of Hopper-specific instruction sets within the implementation, as the codebase primarily targets Ampere architectures for Tensor Cores. Research projects such as ThunkerKitten~\citep{spector2024thunder} and cuDNN 9~\citep{cudnn9} have demonstrated that incorporating Hopper-optimized instructions and tiling abstractions enables both accelerated attention calculation and streamlined implementation processes.

At a more foundational level, the algorithm underlying \textsc{FlashAttention-2} operates within a streamlined, synchronous paradigm, without incorporating asynchrony or low-precision computation into its architecture. Asynchrony, in modern hardware, emerges from the use of specialized components that accelerate essential operations within ML workloads: for instance, matrix multiplications executed by Tensor Cores, or memory operations offloaded to the Tensor Memory Accelerator (TMA), both distinct from general-purpose CUDA cores responsible for arithmetic and control logic. The adoption of lower-precision formats, such as FP8 (introduced with Hopper) or FP4 (in Blackwell), extends the established trajectory seen with FP16 (Pascal, 2017) and BF16 (Ampere, 2020); this approach consistently delivers two to four times greater throughput at equivalent power and silicon area. We examine the innovations in this domain offered by Hopper in \cref{subsec:hardware}. 

Adapting \textsc{FlashAttention-2} to harness these hardware advancements presents technical obstacles: exploiting asynchrony necessitates orchestrating computations so that matrix multiplications and softmax operations can be executed in an overlapped manner, despite intrinsic data dependencies. Moreover, leveraging low-precision arithmetic requires careful strategies to control and mitigate quantization artifacts, which is particularly crucial when dealing with outlier features prevalent in LLMs~\citep{dettmers2208llm, sun2024massive}.

For this purpose, we introduce \textsc{FlashAttention-3}, integrating three novel concepts aimed at enhancing efficiency on contemporary GPU architectures.\footnote{Our discussion is situated within the framework of NVIDIA's Hopper architecture, though the algorithm is compatible with any GPU architecture that offers strong asynchronous processing and supports low-precision computation.}

\begin{enumerate}

\item \textbf{Producer-Consumer asynchrony:} We introduce a warp-specialized approach to software pipelining that leverages the asynchronous execution characteristics of data transfer operations and Tensor Cores. By distributing the roles of data producers and consumers across distinct warps, our scheme enhances the ability of the algorithm to mask the latencies associated with memory accesses and instruction dispatch.
\item \textbf{Hiding softmax under asynchronous block-wise GEMMs:} Our approach interleaves the execution of softmax-related non-GEMM computations, such as floating point multiply-add and exponentiation—operations typically characterized by lower throughput—with asynchronously launched WGMMA instructions responsible for GEMM computation. To facilitate this overlap, we revise the \textsc{FlashAttention-2} algorithm to break existing sequential dependencies between softmax and GEMM steps. As an illustration, in the two-stage implementation of our method, one block of the scores matrix is processed by the softmax routine simultaneously as WGMMA processes the subsequent block asynchronously.
\item \textbf{Hardware-accelerated low-precision GEMM:} The forward pass is restructured to exploit FP8 Tensor Cores for GEMM, which nearly doubles observed TFLOPs/s. To address the disparate data layout expectations between FP32 accumulators and FP8 operand matrices required by WGMMA, we employ block quantization and incoherent processing strategies, thus mitigating the precision loss inherent in transitioning to FP8 computation.

To assess the effectiveness of our approach, we conduct empirical benchmarks of \textsc{FlashAttention-3} on an H100 SXM5 GPU across diverse parameter settings. Our findings indicate that (1) for the forward pass, FP16 achieves between 1.5$\times$ and 2.0$\times$ acceleration relative to \textsc{FlashAttention-2} (achieving up to 740 TFLOPs/s) and delivers 1.5-1.75$\times$ speedup in the backward pass; (2) FP8 execution approaches performance levels of 1.2 PFLOPs/s; and (3) with increasing sequence length, FP16 consistently outperforms alternatives, while FP8 provides competitive throughput compared to the leading attention kernel in NVIDIA's cuDNN library.\footnote{Specifically, for a head dimension of 64, \textsc{FlashAttention-3} FP8 leads; for head dimensions of 128 and 256, it matches for cases without causal masking, though it lags when causal masking is applied.} We further confirm that FP16 \textsc{FlashAttention-3} preserves the same level of numerical error as \textsc{FlashAttention-2}, and performs better than the conventional attention computation—owing to retention of intermediate values (such as softmax rescaling) in FP32. In addition, FP8 \textsc{FlashAttention-3} with block quantization and incoherent processing achieves accuracy 2.6$\times$ greater than standard attention applying per-tensor quantization, especially in scenarios involving outlier features.

\textsc{FlashAttention-3} has been released under a permissive open-source license\footnote{\textsc{FlashAttention-3}
  is available at \texttt{https://github.com/Dao-AILab/flash-attention}}. We aim to make it accessible to a broad research and development community by working toward integration with PyTorch and Hugging Face libraries.

\section{Background: Multi-Head Attention and GPU Characteristics}
\label{sec:background}

\subsection{Multi-Head Attention}
\label{subsec:multi_head_attn}

Let $\mathbf{Q}, \mathbf{K}, \mathbf{V} \in \mathbb{R}^{N \times d}$ be the query, key and value input sequences associated to a single head, where $N$ is the sequence length and $d$ is the head dimension. Then the attention output $\mathbf{O}$ is computed as:

\begin{equation*}
  \mathbf{S} = \alpha \mathbf{Q} \mathbf{K}^\top \in \mathbb{R}^{N \times N}, \quad \mathbf{P} = \mathrm{softmax}(\mathbf{S}) \in \mathbb{R}^{N \times N}, \quad \mathbf{O} = \mathbf{P}\mathbf{V} \in \mathbb{R}^{N \times d},
\end{equation*}

In this context, the $\mathrm{softmax}$ function is executed across each row, and the scaling parameter $\alpha$ is conventionally chosen as $1/\sqrt{d}$. To address potential numerical instability associated with the exponential calculation, it is common practice to subtract $\mathrm{rowmax}(\mathbf{S})$ from $\mathbf{S}$. Within the multi-head attention (MHA) mechanism, every attention head employs distinct projections for queries, keys, and values. These computations are performed in parallel across all heads and batches, yielding the complete output tensor.

Now let $\phi$ be a scalar loss function and let $\mathbf{d}(-) = \partial \phi / \partial (-)$ be notation for the gradient. Given the output gradient $\mathbf{dO} \in \mathbb{R}^{N \times d}$, we compute $\mathbf{dQ}$, $\mathbf{dK}$, and $\mathbf{dV}$ according to the chain rule as follows:

\begin{align*}
  \mathbf{dV} &= \mathbf{P}^\top \mathbf{dO} \in \mathbb{R}^{N \times d} \\
  \mathbf{dP} &= \mathbf{dO} \mathbf{V}^\top \in \mathbb{R}^{N \times N} \\
  \mathbf{dS} &= \mathrm{dsoftmax} (\mathbf{dP}) \in \mathbb{R}^{N \times N} \\
  \mathbf{dQ} &= \alpha \mathbf{dS} \mathbf{K} \in \mathbb{R}^{N \times d} \\
  \mathbf{dK} &= \alpha \mathbf{dS}^\top \mathbf{Q} \in \mathbb{R}^{N \times d},
\end{align*}

In this context, the relationship is expressed as $\mathbf{d}s = (\mathrm{diag}(p) - p p^\top)\mathbf{d}p$, where $p = \mathrm{softmax}(s)$ is parameterized by the vector $s$. We denote the application of this expression to each row individually by $\mathrm{dsoftmax}(\mathbf{dP})$. Moreover, this operation is amenable to parallel execution over both the head and batch dimensions during the backward propagation stage of MHA.

\subsection{GPU hardware characteristics and execution model}
\label{subsec:hardware}

We outline key features of the GPU execution model pertinent to \textsc{FlashAttention-3}, emphasizing the NVIDIA Hopper architecture as a specific example of how this model is implemented.

\paragraph{Memory hierarchy:} The organization of GPU memory follows a hierarchical structure, where higher-bandwidth memories typically have smaller capacities (\cref{tab:gpu-hierarchy})\footnote{\citet{luo2024benchmarking} indicates that shared memory offers a bandwidth of 128 bytes per clock cycle per SM; to estimate the total, we use 132 SMs and a boost clock frequency of 1830 MHz.}. At the highest level, global memory (GMEM)—or HBM—serves as off-chip DRAM accessible by all streaming multiprocessors (SMs). This GMEM content is automatically cached into a chip-resident L2 cache. Further down, each SM features its own highly banked, on-chip cache known as shared memory (SMEM), which is under direct programmer control. At the lowest level within each SM lies the register file.

\paragraph{Thread hierarchy:} In the GPU programming paradigm, execution entities are systematically arranged into a hierarchical structure. At the most granular level are individual threads, which collectively form warps (consisting of 32 threads each). Groups of four adjacent warps make up a warpgroup. These are then organized into threadblocks, also known as cooperative thread arrays (CTAs), which can be further grouped into threadblock clusters in Hopper architectures. At the broadest level, the collection of all threadblocks forms a grid.

The two hierarchies exhibit a strong interdependence. Threads belonging to an individual CTA are scheduled together on the same SM, while CTAs within a single cluster are assigned collectively to the same GPC. All threads in a CTA share direct access to SMEM, whereas each individual thread is allocated no more than 256 private registers (RMEM).

\paragraph{Asynchrony and warp-specialization:}

GPUs achieve high throughput by leveraging parallelism and asynchronous operations to mask both memory and computational delays. To enable asynchronous data transfers between GMEM and SMEM, Hopper introduces a specialized hardware component known as the Tensor Memory Accelerator (TMA) \cite[\S7.29]{cuda}. Moreover, in contrast to preceding generations like Ampere, Hopper's Tensor Core, accessible through the warpgroup-level WGMMA instruction \cite[\S9.7.14]{ptx}, operates asynchronously and is capable of directly obtaining its inputs from shared memory.

The introduction of hardware mechanisms for asynchrony enables the creation of warp-specialized kernels, in which the warps within a CTA are organized strictly as either producers or consumers, restricting their responsibilities solely to data transfers or computation, respectively. This arrangement typically enhances the compiler’s ability to produce efficient instruction orderings \citep{warp-specialization-2011}. Furthermore, Hopper includes functionality for dynamically redistributing registers among different warpgroups using \verb|setmaxnreg| \cite[\S9.7.17.1]{ptx}. This allows warps engaged in MMAs to be allocated more RMEM resources compared to those merely performing TMA operations, which typically require only a single active thread.

\paragraph{Low-precision number formats:}
\label{sec:low-precision-gpu}
Contemporary GPUs incorporate dedicated hardware modules designed to enhance the efficiency of low-precision arithmetic. For instance, utilizing the WGMMA instruction on Hopper enables computation on FP8 Tensor Cores, resulting in a per-SM throughput that is twice as high as that achieved with FP16 or BF16.

Nonetheless, utilizing FP8 WGMMA correctly requires familiarity with the operand layout restrictions. Consider a GEMM operation that multiplies $A \times B^{\top}$, where $A$ has dimensions $M\times K$ and $B$ has dimensions $N\times K$. We define an operand (either $A$ or $B$) as \emph{mn-major} if its outer $M$ or $N$ dimension is contiguous, and as \emph{k-major} if contiguity is found along the inner $K$ dimension. While FP16 WGMMA can accept input operands in either mn-major or k-major configuration when stored in SMEM, FP8 WGMMA restricts operand input to the k-major layout exclusively. Additionally, in contexts like attention mechanisms where one seeks to combine multiple GEMM operations within a single kernel, mismatches between the layout of the FP32 accumulator and FP8 operands hinder the seamless chaining of FP8 WGMMAs.

When considering attention mechanisms, these layout constraints necessitate specific adjustments to the FP8 algorithm’s design, as detailed in \cref{sec:algofp8}.

\subsection{Standard Attention and Flash Attention} In line with the terminology introduced by \citet{dao2022flashattention}, we use the term \textbf{standard attention} to refer to a GPU-based attention implementation that stores the intermediate matrices $\mathbf{S}$ and $\mathbf{P}$ in HBM. The core contribution of \textsc{FlashAttention} is the introduction of a localized softmax reduction that minimizes costly intermediate memory accesses, enabling the fusion of the attention computation within a single GPU kernel. This localized softmax operation is implemented in lines \ref{code-ws:softmax_start}-\ref{code-ws:softmax_end} of the consumer mainloop in \cref{alg:flash3_wgmma_ws_only}, and incorporates the necessary rescaling of blocks in $\mathbf{O}$. A straightforward proof that this approach correctly produces $\mathbf{O}$ is provided in \cite[\S 2.3.1]{dao2023flashattention2}.

\section{FlashAttention-3: Algorithm}
\label{sec:algo}

This section provides an overview of the \textsc{FlashAttention-3} algorithm. We primarily discuss the forward computation here; details regarding the backward pass are given in~\cref{sec:algo_ws_bwd}. Initially, we outline the incorporation of warp-specialization, together with a circular SMEM buffer, within the framework established by \textsc{FlashAttention-2}. Subsequently, we elaborate on leveraging the asynchronous characteristics of WGMMA to develop a two-stage pipeline that overlaps the GEMM and softmax operations. Lastly, we detail the necessary adjustments for FP8, addressing both layout compatibility and the maintenance of accuracy through block quantization and incoherent processing strategies.

\subsection{Producer-Consumer asynchrony through warp-specialization and pingpong scheduling}
\label{sec:algo_ws}

\paragraph{Warp-specialization}
Similar to the approach taken in \textsc{FlashAttention-2}, the forward computation in \textsc{FlashAttention-3} exhibits trivial parallelism across the batch dimension, number of attention heads, and the query sequence length. Therefore, a description at the CTA (Cooperative Thread Array) granularity, focused on one query matrix tile $\mathbf{Q}_i$ generating the corresponding output tile $\mathbf{O}_i$, is sufficient. For clarity, the following discussion outlines the warp-specialization framework using a circular SMEM buffer, but without addressing overlap between the GEMM and softmax computations. Let $d$ denote the size of each attention head, $N$ the total sequence length, and select a query block size $B_r$; this partitions $\mathbf{Q}$ into $T_r = \lceil \frac{N}{B_r} \rceil$ submatrices $\mathbf{Q}_1, .., \mathbf{Q}_{T_r}$.

\begin{algorithm}[H]
    \caption{\small\label{alg:flash3_wgmma_ws_only}\textsc{FlashAttention-3} forward pass \textbf{excluding} intra-consumer overlap -- CTA perspective}
    \begin{algorithmic}[1]
\REQUIRE Input matrices $\mathbf{Q}_i \in \mathbb{R}^{B_r \times d}$, $\mathbf{K}, \mathbf{V} \in \mathbb{R}^{N \times d}$ stored in HBM; key block size $B_c$, compute $T_c = \lceil \frac{N}{B_c} \rceil$.
\STATE Set up pipeline control for barrier synchronization using an $s$-stage circular SMEM buffer.
\IF {executing as producer warpgroup}
\STATE Free a predetermined quantity of registers.
\STATE Launch data movement of $\mathbf{Q}_i$ from HBM to shared memory.
\STATE Upon data arrival, signal to consumer the successful transfer of $\mathbf{Q}_i$.
\FOR{$0 \le j < T_c$}
    \STATE Stall until the $(j\,\%\,s)$ stage in the buffer is processed.
    \STATE Dispatch loads for $\mathbf{K}_j, \mathbf{V}_j$ from HBM to shared memory within the $(j\,\%\,s)$ buffer stage.
    \STATE After the load, notify consumer groups about the availability of $\mathbf{K}_j, \mathbf{V}_j$.
\ENDFOR
\ELSE \STATE Allocate registers as required by the count of consumer warps.
\STATE Within on-chip memory, set $\mathbf{O}_i = (0) \in \mathbb{R}^{B_r \times d}$, and initialize $\ell_i = 0$ and $m_i = -\infty$ in $\mathbb{R}^{B_r}$.
\STATE Hold execution until $\mathbf{Q}_i$ is resident in shared memory.
\FOR{$0 \le j < T_c$}
\STATE Hold until the buffer holds $\mathbf{K}_j$.
\STATE Calculate $\mathbf{S}_i^{(j)} = \mathbf{Q}_i \mathbf{K}_j^T$ by SS-GEMM, and synchronize as needed. \label{alg:ws_only_gemm1}
\STATE Set $m_i^{\mathrm{old}} = m_i$, then update $m_i = \max(m_i^{\mathrm{old}}, \mathrm{rowmax}(\mathbf{S}_i^{(j)}))$. \label{code-ws:softmax_start}
\STATE Determine $\widetilde{\mathbf{P}}_i^{(j)} = \mathrm{exp}(\mathbf{S}_i^{(j)} - m_i)$, updating $\ell_i = \mathrm{exp}(m_i^{\mathrm{old}} - m_i) \ell_i + \mathrm{rowsum}(\widetilde{\mathbf{P}}_i^{(j)})$. \label{code-ws:softmax_end}
\STATE Wait until $\mathbf{V}_j$ becomes available in shared memory.
\STATE Update $\mathbf{O}_i = \mathrm{diag}(\mathrm{exp}(m_i^{\mathrm{old}} - m_i))^{-1} \mathbf{O}_i + \widetilde{\mathbf{P}}_i^{(j)} \mathbf{V}_j$ utilizing RS-GEMM, then commit and synchronize. \label{alg:ws_only_gemm2}
\STATE Release the $(j\,\%\,s)$ buffer stage for use by the producer.
\ENDFOR
\STATE Post-process $\mathbf{O}_i = \mathrm{diag}(\ell_i)^{-1} \mathbf{O}_i$ and compute $L_i = m_i + \log(\ell_i)$.
\STATE Output $\mathbf{O}_i$ and $L_i$ to HBM in their respective positions in $\mathbf{O}$ and $L$.
\ENDIF
\end{algorithmic}
\end{algorithm}

In our implementation of \cref{alg:flash3_wgmma_ws_only} on Hopper, we utilize \verb|setmaxnreg| to manage register (de)allocations, employ TMA for loading $\mathbf{Q}_i$ and $\{ \mathbf{K}_j, \mathbf{V}_j \}_{0 \leq j < T_c}$, and use WGMMA to perform GEMMs within the consumer mainloop. The SS and RS prefixes designate whether the initial operand resides in shared memory or the register file, respectively. When analyzing the runtime behavior of \cref{alg:flash3_wgmma_ws_only}, it is important to recognize that TMA load instructions benefit from asynchronous operation, and thus, the initiation of a new TMA load is not contingent upon the completion of previous loads. Additionally, during the initial $s$ iterations of the producer mainloop, waits are omitted as the buffer has not yet reached full occupancy.

\paragraph{Pingpong scheduling} The asynchronous execution of WGMMA and TMA, together with warp-specialization, makes it possible to interleave the softmax calculation from one warpgroup while a different warpgroup is performing GEMM. This is particularly advantageous given that, on modern computing hardware, operations other than matrix multiplication are executed at significantly lower throughput compared to matmul. For illustration, consider the H100 SXM5 GPU, which can achieve up to 989 TFLOPS for FP16 matmul operations, whereas the throughput for special functions like the exponential\footnote{The CUDA programming guide states that each streaming multiprocessor (SM) is capable of performing 16 special function operations per clock cycle. By multiplying 16 operations by 132 SMs and a clock frequency of 1830 MHz, we derive a total of 3.9 TFLOPS for special functions.} (an essential function for softmax) is only 3.9 TFLOPS. In the case of an FP16 attention forward pass with a head size of 128, the volume of matmul FLOPS exceeds that of exponential operations by a factor of 512; however, since the exponential’s achievable throughput is 256 times smaller, the time required by the exponential function can occupy up to 50\% of the total cycle time relative to matmul. This discrepancy becomes more pronounced when switching to FP8: although the matmul rate doubles, the throughput for exponential computations does not increase, amplifying the imbalance.

Because the exponential function is executed by a distinct hardware component (the multi-function unit), optimal performance requires aligning the exponential computations with matrix multiplication operations executed by the Tensor Cores. To achieve this coordination, we implement synchronization barriers (\texttt{bar.sync} instructions), which enforce an execution order: GEMMs (specifically, GEMM1 -- $\mathbf{P} \mathbf{V}$ from the current iteration and GEMM0 -- $\mathbf{Q} \mathbf{K}^\top$ from the subsequent iteration) associated with warpgroup 1 are executed before those associated with warpgroup 2. Consequently, while warpgroup 1 computes softmax, warpgroup 2 proceeds with its GEMMs. This process then reverses, allowing warpgroup 2 to execute softmax while warpgroup 1 returns to its GEMM computations, thus establishing a “pingpong” scheduling scheme. This behavior is depicted in~\cref{fig:pingpong_scheduling}. 

Although real-world execution of this pingpong approach exhibits deviations from the idealized scheduling illustrated in the figure, we observe notable performance gains in practice—for example, increasing throughput from 570 TFLOPS to 620-640 TFLOPS on FP16 forward passes with a head dimension of 128 and sequence length of 8192.

\paragraph{Attention variants} In the case of multi-query attention~\citep{shazeer2019fast} and grouped query attention~\citep{ainslie2023gqa}, our method mirrors that of \textsc{FlashAttention-2} by modifying tensor indexing procedures so as to prevent redundant copies of $\mathbf{K}$ and $\mathbf{V}$ within HBM.

\subsection{Intra-warpgroup overlapping GEMMs and softmax}

It is possible to interleave certain softmax instructions with GEMM instructions, even within a single warpgroup. Here, we outline a specific method for achieving this overlap.

In the attention algorithm, operations within the inner loop (main loop) have sequential dependencies that impede parallelization within a single iteration. For example, (local) softmax (lines \ref{code-ws:softmax_start} to \ref{code-ws:softmax_end}) relies on the output $\mathbf{S}_i^{(j)}$ of the first GEMM, while the second GEMM takes its result $\widetilde{\mathbf{P}}_i^{(j)}$ as an operand. Indeed, the wait statements in lines \ref{alg:ws_only_gemm1} and \ref{alg:ws_only_gemm2} of \cref{alg:flash3_wgmma_ws_only} serialize the execution of softmax and GEMMs. However, we can break these dependencies by pipelining across iterations through additional buffers in registers. Pursuing this idea, we propose the following two-stage\footnote{Note that the number of stages of the overlapping scheme is bounded by, but need not equal, the number $s$ of stages in the circular SMEM buffer.} GEMM-softmax pipelining algorithm:

\begin{algorithm}[H]
    \caption{\small\label{alg:flash3_wgmma}\textsc{FlashAttention-3} consumer warpgroup forward pass}
    \begin{algorithmic}[1]
      \REQUIRE  Input matrices $\mathbf{Q}_i \in \mathbb{R}^{B_r \times d}$ and $\mathbf{K}, \mathbf{V} \in \mathbb{R}^{N \times d}$ located in HBM, a specified key block size $B_c$, and $T_c = \lceil \frac{N}{B_c} \rceil$.
      \STATE Adjust the allocation of registers according to the total number of consumer warps.
      \STATE Set $\mathbf{O}_i = (0) \in \mathbb{R}^{B_r \times d}$, $\ell_i = (0) \in \mathbb{R}^{B_r}$, and $m_i = (-\infty) \in \mathbb{R}^{B_r}$ in on-chip memory.
      \STATE Suspend further operations until both $\mathbf{Q}_i$ and $\mathbf{K}_0$ are fully loaded into shared memory.
      \STATE Utilize WGMMA to calculate $\mathbf{S}_{\mathrm{cur}} = \mathbf{Q}_i \mathbf{K}_0^T$; issue commit and await completion.
      \STATE Make the $0$th stage of the $\mathbf{K}$ buffer available for future use.
      \STATE Determine $m_{i}$, update $\tilde{\mathbf{P}}_{\mathrm{cur}}$ and $\ell_{i}$ from $\mathbf{S}_{\mathrm{cur}}$, and perform a rescaling of $\mathbf{O}_i$ accordingly.
      \FOR{$1 \le j < T_c - 1$}
        \STATE Pause execution until $\mathbf{K}_j$ has been copied to shared memory.
        \label{code:mainloop_start}
        \STATE Initiate calculation of $\mathbf{S}_{\mathrm{next}} = \mathbf{Q}_i \mathbf{K}_{j}^T$ via WGMMA; commit the operation without synchronizing immediately. \label{code:first_wgmma}
        \STATE Wait for the transfer of $\mathbf{V}_{j-1}$ into shared memory to finish.
        \STATE Perform the computation $\mathbf{O}_{i} = \mathbf{O}_{i} + \tilde{\mathbf{P}}_{\mathrm{cur}} \mathbf{V}_{j-1}$ using WGMMA; commit the task while deferring synchronization. \label{code:second_wgmma}
        \STATE Ensure that the previous WGMMA operation $\mathbf{Q}_i \mathbf{K}_{j}^T$ has completed before proceeding.
        \STATE Based on the values in $\mathbf{S}_{\mathrm{next}}$, update $m_{i}$, calculate $\tilde{\mathbf{P}}_{\mathrm{next}}$ and revise $\ell_{i}$. \label{code:softmax}
        \STATE Synchronize once the $\tilde{\mathbf{P}}_{\mathrm{cur}} \mathbf{V}_{j-1}$ WGMMA operation finalizes, then reapply scaling to $\mathbf{O}_i$.
        \STATE Free the buffer stage indexed by $(j\,\%\,s)$ for $\mathbf{K}$, and by $(j-1\,\%\,s)$ for $\mathbf{V}$.
        \STATE Transfer contents from $\mathbf{S}_{\mathrm{next}}$ to $\mathbf{S}_{\mathrm{cur}}$ for the upcoming iteration.
        \label{code:mainloop_end}
      \ENDFOR \STATE Wait until $\mathbf{V}_{T_c - 1}$ resides in shared memory.
      \STATE Execute the update
      $\mathbf{O}_{i} = \mathbf{O}_{i} + \tilde{\mathbf{P}}_{\mathrm{last}} \mathbf{V}_{T_c - 1}$ using WGMMA, and finalize with a commit and wait.

\STATE Epilogue: Adjust the scale of $\mathbf{O}_{i}$ according to $m_{i}$. Calculate $L_{i}$ utilizing both $m_{i}$ and $\ell_{i}$.
      Store $\mathbf{O}_{i}$ and $L_{i}$ into HBM in the $i$-th positions of $\mathbf{O}$ and $L$, respectively.
    \end{algorithmic}
  \end{algorithm}

\cref{alg:flash3_wgmma} serves as a substitute for the consumer component found in \cref{alg:flash3_wgmma_ws_only}, thereby realizing the complete \textsc{FlashAttention-3} procedure for FP16 computations. Broadly, the use of WGMMA in this context denotes asynchronous GEMM. Inside the mainloop (spanning lines \ref{code:mainloop_start} through \ref{code:mainloop_end}), the execution of the second WGMMA instruction during iteration $j$ (as shown in line \ref{code:second_wgmma}) is executed concurrently with the softmax calculations of the subsequent iteration $j+1$ (see line \ref{code:softmax}).

Although the pipelined architecture discussed above has the potential to yield significant performance improvements in theory, several real-world considerations may impact its effectiveness:
\paragraph{Compiler reordering} The pseudocode outlines a best-case scenario in terms of execution sequence; however, in practice, the NVCC compiler frequently reorders instructions to enhance performance. Such automatic instruction reorganization can disturb the deliberate interleaving of WGMMA and non-WGMMA operations, possibly reducing the intended advantages of the pipeline or introducing anomalies in execution. Examination of the generated SASS code confirms that, in this study, the compiler does create the expected overlapped execution (\cref{sec:2-stage-sass}).
\paragraph{Register pressure} To fully exploit the potential of the pipeline, register spilling must be kept to a minimum. Nonetheless, the two-stage pipelined approach inherently increases the register requirement, as it must buffer intermediate values and manage execution context between stages. More specifically, an additional buffer for $\mathbf{S}_{\mathrm{next}}$ must reside in registers, resulting in an overhead of $B_r \times B_c \times \text{sizeof}(\text{float})$ registers per threadblock. This elevated register footprint may compete with efforts to deploy larger block sizes—another technique aimed at performance enhancement that also demands significant register resources. Hence, practical tuning should involve careful profiling to balance these competing needs.
\paragraph{3-stage pipelining} Building on the foundation of the 2-stage approach, we introduce a 3-stage scheme designed to further overlap the second WGMMA step with the softmax computation. This added pipeline depth can potentially maximize the throughput of Tensor Cores, but comes at the cost of requiring even more registers to store the additional intermediate results. Consequently, finding an optimal balance between the tile dimensions and the pipeline depth becomes increasingly complex. Details on the implementation and evaluation of this 3-stage pipelined algorithm are presented in ~\cref{sec:3-stage}.

\subsection{Low-precision with FP8}
\label{sec:algofp8}

\textbf{Efficiency: layout transformations.} Executing the forward propagation of \textsc{FlashAttention-3} using FP8 precision introduces extra difficulties regarding layout compatibility that do not arise with FP16 precision.

To begin, it is important to highlight that the input tensors $\mathbf{Q}$, $\mathbf{K}$, and $\mathbf{V}$ are usually arranged contiguously along the head dimension. However, the FP8 WGMMA’s k-major requirement for the second GEMM dictates that $\mathbf{V}$—specifically, the portions of $\mathbf{V}$ brought into SMEM—must be contiguous with respect to the sequence length dimension. Since the TMA load operation is unable to alter which dimension is contiguous, this necessitates one of two approaches: (1) performing a pre-processing transpose of $\mathbf{V}$ within GMEM, or (2) applying an in-kernel transpose to portions of $\mathbf{V}$ after these segments have been loaded into SMEM. For strategy (1), the transpose can be either (1a) integrated with the epilogue of an earlier step like the rotary embedding, or (1b) executed as an independent pre-processing transpose kernel\footnote{A high-performance transpose kernel is capable of reaching near-peak device bandwidth~\citep{colfax_cutlass_transpose_2024}.} to swap the head and sequence length strides. Yet, option (1a) presents challenges when attempting to incorporate it into general-purpose libraries, and option (1b) introduces significant inefficiency—particularly under the memory-bound conditions typical of inference workloads.

For FP8 \textsc{FlashAttention-3}, we select option (2) as our approach. Within the kernel, we utilize LDSM (\verb|ldmatrix|) and STSM (\verb|stmatrix|) instructions, which facilitate cooperative loading from SMEM to RMEM and storing from RMEM to SMEM by a thread warp in 128-byte segments.\footnote{Although the PTX documentation characterizes LDSM and STSM as operations on $8 \times 8$ matrices with 16-bit elements \cite[\S 9.7.13.4.15-16]{ptx}, in our implementation with FP8, we combine pairs of 8-bit elements to fit within the constraints of these instructions. Still, the transposing variants of LDSM/STSM cannot directly separate packed 8-bit items, making it necessary to perform some intermediary register manipulations between LDSM and STSM to achieve a tile-wise transpose; we do not elaborate these specifics here.} The LDSM/STSM instructions not only provide good register efficiency—permitting their use within the producer warpgroup—but also support layout transpositions during data transfers. In addition, from the second iteration onward, the transpose of the subsequent $\mathbf{V}$ tile can be scheduled to overlap with the computation of two WGMMAs that process the prior $\mathbf{V}$ tile along with the current $\mathbf{K}$ tile.

Second, we note that, in contrast to FP16, the FP32 accumulator employed by FP8 WGMMA operations possesses a memory layout distinct from that anticipated for operand A when this operand resides in registers. \cref{fig:rmem_accum} and \cref{fig:rmem_operand} illustrate segments of both memory layouts, where each entry is mapped to thread-held registers following the shown sequence. By leveraging byte permute instructions, the accumulator output from the initial WGMMA can be reorganized to align with the memory arrangement required by a subsequent WGMMA, also corresponding with the format of the $\mathbf{V}$ tile after an in-kernel transpose. Concretely, as depicted in \cref{fig:rmem_accum}, the byte ordering is transformed in the pattern
$$\{ \verb|d0 d1 d4 d5 d2 d3 d6 d7| \},$$
with this register rearrangement repeated for every 8-byte segment. Logically, for the $\mathbf{P}$ tile, this operation equates to a permutation of its columns, such that, for instance, columns $0189$ are brought to the foremost four columns. Ensuring that WGMMA subsequently generates the intended output tile then reduces to coordinating the in-kernel transpose to produce a corresponding row rearrangement for the $\mathbf{V}$ tile.\footnote{The ability to implement this transformation within the in-kernel transpose removes the necessity of utilizing shuffle instructions to reassign register contents among threads, a process discussed previously in~\citep{colfax_fp8_flashattention_2024}.}

\textbf{Accuracy: block quantization and incoherent processing.} The FP8 (e4m3) format allocates 3 bits for the mantissa and 4 bits for the exponent, which leads to increased numerical errors compared to FP16 or BF16 formats. Furthermore, in large-scale models, it is common to encounter outlier values~\citep{dettmers2208llm, sun2024massive} whose magnitudes far exceed the majority of tensor entries, making the process of quantization more challenging. A standard approach is to apply per-tensor scaling~\citep{micikevicius2022fp8}, in which a separate scaling factor is stored for each tensor (such as one for $\mathbf{Q}$, one for $\mathbf{K}$, and one for $\mathbf{V}$). 

To address the heightened numerical error associated with attention in FP8, we utilize two primary strategies:

\begin{enumerate}

\item \textbf{Block quantization}: In this strategy, each block retains a single scalar, meaning that for tensors $\mathbf{Q}$, $\mathbf{K}$, and $\mathbf{V}$, we partition them into blocks with dimensions $B_r \times d$ or $B_c \times d$ and apply quantization independently to each block. This blockwise quantization process can be seamlessly integrated with preceding operations (such as rotary embedding) without incurring additional latency, given that rotary embeddings are constrained by memory bandwidth. Moreover, since the \textsc{FlashAttention-3} algorithm operates on blockwise structures, it allows for the per-block scaling of
$\mathbf{S}$ to reflect the applied quantization, without additional computational overhead.

\item \textbf{Incoherent processing}: To mitigate the influence of outlier values, $\mathbf{Q}$ and $\mathbf{K}$ are premultiplied by a random orthogonal matrix $\mathbf{M}$ prior to FP8 quantization. Owing to the orthogonality property of $\mathbf{M}$, i.e., $\mathbf{M} \mathbf{M}^\top = I$, we have $(\mathbf{Q} \mathbf{M})(\mathbf{K} \mathbf{M})^\top = \mathbf{Q} \mathbf{K}^\top$, thus ensuring that this transformation leaves the attention calculation unchanged. This approach effectively disperses outliers, since each component of $\mathbf{Q} \mathbf{M}$ or $\mathbf{K} \mathbf{M}$ forms a random linear combination of $\mathbf{Q}$ or $\mathbf{K}$ elements, lessening the quantization error. Following the methodology of \citet{chee2024quip} and~\citet{tseng2024quip}, we construct $\mathbf{M}$ as the product of random diagonal matrices with entries $\pm 1$ and a Hadamard matrix. This structure allows for multiplication to be performed in $O(d \log d)$ time, compared to the naive $O(d^2)$, and this step can also be fused with the rotary embedding without incurring extra computation.

\end{enumerate}

We demonstrate empirically in \cref{sec:numerical_error} that these two methods reduce numerical errors by up to a factor of 2.6.

\section{Empirical Validation}
\label{sec:exp}

To implement \textsc{FlashAttention-3} and assess both its performance and precision, we leverage CUTLASS~\citep{Thakkar_CUTLASS_2023} primitives, specifically the WGMMA and TMA abstractions.

\begin{itemize}

\item \textbf{Attention Benchmarking.} We evaluate the performance of \textsc{FlashAttention-3} by recording its execution times over varying sequence lengths and contrasting the results with those from PyTorch's reference implementation, \textsc{FlashAttention-2}, a Triton-based version of \textsc{FlashAttention-2} leveraging H100-targeted instructions, and a vendor-optimized \textsc{FlashAttention-2} implementation in cuDNN for H100 GPUs. Our findings demonstrate that \textsc{FlashAttention-3} achieves speedups of up to 2.0$\times$ over \textsc{FlashAttention-2} and up to 1.5$\times$ over its Triton counterpart. Moreover, \textsc{FlashAttention-3} attains a throughput of up to 740 TFLOPs/s, corresponding to 75\% of H100 GPUs' peak theoretical performance.
\item \textbf{Ablation Analysis.} We verify that the observed acceleration in \textsc{FlashAttention-3} can be attributed to specific innovations in our methodology, particularly warp-specialization and the GEMM-softmax pipelining technique.
\item \textbf{FP8 Attention Accuracy.} Our experiments show that employing block quantization and incoherent processing effectively reduces the numerical error in FP8 \textsc{FlashAttention-3} by a factor of 2.6$\times$.

\subsection{Benchmarking Attention}
\label{subsec:benchmark_attn}

We evaluate the execution time of various attention algorithms on an H100 80GB SXM5 GPU under different configurations (with and without a causal mask, using head dimensions of 64 or 128) and FP16 data precision. The performance outcomes are presented in~\cref{fig:benchmark_attn_fwd} and~\cref{fig:benchmark_attn_bwd}. Our findings indicate that \textsc{FlashAttention-3} achieves a speedup of approximately 1.5-2.0$\times$ over \textsc{FlashAttention-2} in the forward computation, and 1.5-1.75$\times$ in the backward computation. In comparison to the conventional attention mechanism, \textsc{FlashAttention-3} demonstrates acceleration by a factor of up to 3-16$\times$. Remarkably, for sequence lengths of 1k and beyond, \textsc{FlashAttention-3} exceeds the throughput of a proprietary library (cuDNN -- closed source) specifically optimized for H100 hardware.

\paragraph{Benchmark settings:} We vary the sequence length as 512, 1k, ..., 16k, and set batch size so that the total number of tokens is 16k. We set the hidden dimension to 2048, and head dimension to be either 64, 128, or 256 (i.e., 32 heads, 16 heads, or 8 heads). To calculate the FLOPs of the forward pass, we use:

\begin{equation*}
  4 \cdot \text{seqlen}^2 \cdot \text{head dimension} \cdot \text{number of heads}.
\end{equation*}

By applying causal masking, we halve this value to reflect that roughly 50% of the computations are performed. For estimating the FLOPs during the backward pass, we take the FLOPs from the forward computation and multiply by 2.5, because the backward pass involves five matrix multiplications while the forward pass only involves two, primarily as a consequence of recomputation.

Additionally, we evaluate the forward pass runtime using FP8 under comparable conditions. The outcomes for headdim 256 are displayed in ~\cref{fig:benchmark_attn_fp8}, with comprehensive results presented in \cref{sec:benchmark_attn_fp8_full}.

\subsection{Ablation Study: 2-Stage Pipelining Experiments}
\label{sec:2-stage-pipelining-experiments}

We conduct an ablation of both the 2-stage WGMMA-softmax pipelining and the warp-specialization techniques on non-causal FP16 \textsc{FlashAttention-3}, keeping the parameters $\{\text{batch}, \text{seqlen}, \text{nheads}, \text{hdim}\} = \{4, 8448, 16, 128\}$ fixed.
As shown in~\cref{table:ablation_pipelining}, these algorithmic enhancements—namely introducing asynchrony through warp-specialization and enabling overlap between GEMM and softmax—yield a considerable performance boost, increasing throughput from 570 to 661 TFLOPs.

\subsection{Numerical Error Validation}
\label{sec:numerical_error}

Given the recent focus on the numerical error of \textsc{FlashAttention}~\citep{golden2024flash}, we perform a comparative analysis involving \textsc{FlashAttention-2}, \textsc{FlashAttention-3}, and a baseline attention implementation, utilizing a reference FP64 implementation for evaluation. To mirror the presence of outlier features and activations commonly observed in LLMs~\citep{dettmers2208llm, sun2024massive}, we sample the entries of $\mathbf{Q}, \mathbf{K}, \mathbf{V}$ according to the following distribution:

\begin{equation*}
  \mathcal{N}(0, 1) + \mathcal{N}(0, 100) \cdot \mathrm{Bernoulli}(0.001).
\end{equation*}

In other words, while most entries are sampled from a normal distribution with mean zero and unit standard deviation, an additional independent noise term with standard deviation 10 is incorporated into 0.1\% of the entries. The root mean squared error (RMSE) for this setup is reported in~\cref{table:numerical_error}. Utilizing FP16, both \textsc{FlashAttention-2} and \textsc{FlashAttention-3} produce RMSE values that are 1.7$\times$ smaller than those from the conventional implementation, due to storing intermediate (softmax) computations in FP32. The FP8 baseline attention employs per-tensor scaling, calculates matrix multiplication accumulators in FP32, and holds softmax intermediates in FP16. Owing to block quantization and its processing method, \textsc{FlashAttention-3} in FP8 achieves 2.6$\times$ the accuracy of this baseline.

\section{Dicussion, Limitations, Conclusion}
\label{sec:discussion}

Through \textsc{FlashAttention-3}, we showcase that leveraging advanced programming strategies and exploiting recent hardware capabilities—particularly asynchrony and low-precision arithmetic—can substantially enhance both the performance and precision of attention mechanisms. Our approach achieves a 1.5-2.0$\times$ speedup over \textsc{FlashAttention-2}, while simultaneously lowering FP8 numerical error by a factor of 2.6 compared to conventional per-tensor quantization schemes. Notwithstanding these improvements, there remain aspects requiring further refinement, such as optimization for LLM inference workloads, implementing a persistent kernel structure within the FP8 kernel,\footnote{In the benchmarks conducted, only the FP16 \textsc{FlashAttention-3} variant employs a persistent kernel and load balancing mechanism; the FP8 \textsc{FlashAttention-3} does not. This partially accounts for the inferior performance of FP8 \textsc{FlashAttention-3} on tasks with limited sequence lengths or requiring causal masking, as opposed to FP8 cuDNN kernels.} and investigating the consequences of employing low-precision attention during large-scale model training. While our evaluation targets Hopper GPUs, we anticipate that the principles and optimizations introduced will extend to a broad range of hardware accelerators. Ultimately, we aspire that making attention both faster and more precise will open doors to new use cases, especially in domains involving extended context processing.

\end{document}